\documentclass[11pt]{article}
% Кодировка и язык
\usepackage[utf8]{inputenc}           % Кодировка исходного файла (UTF-8)
\usepackage[T2A]{fontenc}             % Кодировка шрифтов для кириллицы
\usepackage[english,russian]{babel}   % Локализация: русские названия, переносы

% Поиск в PDF
\usepackage{cmap}                     

% Поля страницы
\usepackage{geometry}
\geometry{left=2cm, right=2cm, top=2cm, bottom=2cm}

% Математика
\usepackage{amsmath, amssymb, amsthm} % Базовые пакеты для математики
\usepackage{mathtools}                % Расширение amsmath (доп. символы и окружения)

% Картинки
\usepackage{graphicx}                 % Вставка картинок (\includegraphics)
\graphicspath{{pictures/}}            % Путь к папке с картинками
\usepackage{float}                    % Позволяет использовать [H] для жёсткого позиционирования
\usepackage{wrapfig}                  % Обтекание текста вокруг картинок

% Таблицы
\usepackage{array}                    % Расширенные возможности для таблиц
\usepackage{booktabs}                 % Красивые таблицы (\toprule, \midrule, \bottomrule)
\usepackage{multirow}                 % Объединение строк в таблицах
\usepackage{tabularx}                 % Таблицы с автоматической шириной колонок
\usepackage{ragged2e}                 % Улучшенное выравнивание (используется в tabularx)

% Списки
\usepackage{enumitem}                 % Гибкая настройка списков (itemize, enumerate)

% Колонтитулы
\usepackage{fancyhdr}                 

% Гиперссылки
\usepackage{xcolor}                   % Работа с цветами
\usepackage{hyperref}                 % Гиперссылки, оглавление, перекрёстные ссылки

% Умные ссылки
\usepackage{cleveref}                 

% Типографика
\usepackage{microtype}                % Улучшает переносы и межсловные расстояния

% Подписи к картинкам и таблицам
\usepackage{caption}                  % Настройка подписей (шрифт, отступы)

% Рисование схем и графиков
\usepackage{tikz}                     % Рисование схем, графиков, диаграмм
\usepackage{pgfplots}                 % Построение графиков функций (на основе tikz)
\pgfplotsset{compat=newest}           % Версия совместимости pgfplots

% Алгоритмы и код
\usepackage{algorithm}                % Окружения для алгоритмов
\usepackage{algorithmic}              % Синтаксис для алгоритмов
\usepackage{listings}                 % Вставка кода с подсветкой синтаксиса


% Настройка нумерации секций и подсекций
\renewcommand{\thesection}{\arabic{section}}                   
\renewcommand{\thesubsection}{\thesection.\arabic{subsection}} 


% Настройка списков: убираем лишние отступы слева
\setlist[itemize]{leftmargin=}
\setlist[enumerate]{leftmargin=}


% Определение окружений для теорем, лемм, определений и т.д.
% Нумерация привязана к номеру секции
\newtheorem{theorem}{Теорема}[section]
\newtheorem{definition}{Определение}[section]
\newtheorem{remark}{Замечание}[section]
\newtheorem{lemma}{Лемма}[section]
\newtheorem{example}{Пример}[section]
\newtheorem{consequence}{Следствие}[section]
\newtheorem{algorithmm}{Алгоритм}
\newtheorem{question}{Вопрос}[section]
\newtheorem{answer}{Ответ}[section]
\newtheorem{statement}{Утверждение}[section]


% Определение цветов
\definecolor{pblue}{rgb}{0.13,0.13,1}    
\definecolor{pgreen}{rgb}{0,0.5,0}       
\definecolor{pred}{rgb}{0.9,0,0}         
\definecolor{pgrey}{rgb}{0.46,0.45,0.48}

\newcommand{\lectureNumber}{5}
\newcommand{\lectureTitle}{Линейные неоднородные ДУ $n$-ого порядка с постоянными коэффициентами. Квазимногочлены. Метод вариации постоянных}
\newcommand{\courseName}{Элементы дифференциальных уравнений}
\newcommand{\semester}{3 семестр, 2026/2027 уч. год}


\DeclareMathOperator{\Ree}{Re}
\DeclareMathOperator{\Imm}{Im}
\DeclareMathOperator{\Mat}{Mat}
\DeclareMathOperator{\Lip}{Lip}
\DeclareMathOperator{\loc}{loc}

\pagestyle{fancy}
\fancyhf{}
\fancyhead[L]{Чил аут окончен, ребятки, свага накоплена}
\fancyhead[R]{\thepage}
\setlength{\headheight}{15pt}

\hypersetup{
    colorlinks=true,
    linkcolor=blue!70!black,
    citecolor=blue!70!black,
    filecolor=blue!70!black,
    urlcolor=blue!70!black,
}

\begin{document}


\begin{titlepage}
    \centering
    
    {\large Московский физико-технический институт \par}
    {\large Высшая школа программной инженерии \par}
    \vspace{2cm}

    {\Large \textbf{\courseName} \par}
    {\large \semester \par}
    
    \vspace*{\fill}

    {\Huge \textbf{Лекция \lectureNumber} \par}
    {\Huge \textbf{\lectureTitle} \par}
    
    \vspace*{\fill}
    
    {\large Лектор: Саулин Сергей Михайлович \par}
    {\large Автор: \href{https://t.me/ceagest}{Юра} \par}

\end{titlepage}

\pagenumbering{arabic}

\section{Линейные неоднородные ДУ $n$-ого порядка с постоянными коэффициентами}

Теперь будем рассматривать уравнения следующего вида:

\begin{equation} \label{eq:1}
    a_0 y^{(n)} + a_1 y^{(n - 1)} + \ldots + a_n y = F(x), \; \; a_0, a_1, \ldots, a_n \in \mathbb{C}, \; \; a_0 \neq 0
\end{equation}

\begin{equation*}
    F \colon I \to \mathbb{R}, \; \; I \subset \mathbb{R} \text{ --- интервал}
\end{equation*} \leavevmode
\\
Нам также понадобится соответствующее старое доброе однородное:
\begin{equation} \label{eq:2}
    a_0 y^{(n)} + a_1 y^{(n - 1)} + \ldots + a_n y = 0
\end{equation}

\begin{remark}
    Следующие два утверждения позволяют чуть больше понять, как же устроено общее решение (\ref{eq:1}):
    \begin{enumerate}
        \item Пусть $y_1(x)$ --- решение \textbf{соответствующего} однородного уравнения (\ref{eq:2}), а $y_2(x)$ --- решение уравнения (\ref{eq:1}). Тогда $y(x) = y_1(x) + y_2(x)$ --- решение уравнения (\ref{eq:1}).
        \item Пусть $y_1(x), y_2(x)$ --- решения уравнения (\ref{eq:1}). Тогда $y_1(x) - y_2(x)$ --- решение уравнения (\ref{eq:2}).
    \end{enumerate}
\end{remark}

\begin{proof}
    Докажем утверждения по очереди.
    \begin{enumerate}
        \item Так как $y_1(x)$ --- решение (\ref{eq:2}), то 
        \[a_0 y_1^{(n)} + \ldots + a_n y_1 = 0\]
        Так как $y_2(x)$ --- решение (\ref{eq:1}), то 
        \[a_0 y_2^{(n)} + \ldots + a_n y_2 = F(x)\]
        Следовательно,
        \[a_0 (y_1 + y_2)^{(n)} + \ldots + a_n (y_1 + y_2) = F(x)\]
        \item Так как $y_1(x)$ и $y_2(x)$ --- решения (\ref{eq:1}), то 
        \[a_0 y_1^{(n)} + \ldots + a_n y_1 = F(x)\]
        \[a_0 y_2^{(n)} + \ldots + a_n y_2 = F(x)\]
        Следовательно,
        \[a_0 (y_1 - y_2)^{(n)} + \ldots + a_n (y_1 - y_2) = 0\]
    \end{enumerate}
\end{proof} \leavevmode
\\
Из этого замечания сразу следует, что общее решение неоднородного уравнения (\ref{eq:1}) есть сумма частного решения (\ref{eq:1}) и общего решения однородного уравнения (\ref{eq:2}). Таким образом, задача об изучении множества решений (\ref{eq:1}) свелась к задаче о поиске частного решения (\ref{eq:1}). Этим мы и будем заниматься дальше.

\subsection{Квазимногочлены}

Будем обозначать за $\mathbb{R}_n[x]$ пространство многочленов с вещественными коэффициентами от переменной $x$ степени $\leq n$. Также условимся, что $\mathbb{R}[x]$ --- кольцо многочленов над полем $\mathbb{R}$, то есть
\[\mathbb{R}[x] = \bigcup_{n \geq 0} \mathbb{R}_n[x]\]

\begin{definition}
    \textbf{Квазимногочленом} $F$ степени $m$ с показателем $\lambda = \alpha + i \beta$ называется функция вида
    \[F(x) = e^{\alpha x} (p(x) \cos \beta x + q(x) \sin \beta x), \; \; p, q \in \mathbb{R}[x], \; \; \max \{\deg p, \deg q\} = m\]
\end{definition}

\begin{definition}
    Определим \textbf{пространство квазимногочленов} фиксированной степени и фиксированного показателя:
    \[Q_{\lambda, m} = \{r(x) = e^{\alpha x} (p(x) \cos \beta x + q(x) \sin \beta x) \; | \; p, q \in \mathbb{R}_{m - 1}[x]\} \]
\end{definition} \leavevmode
\\
Из определения вытекает очевидное наблюдение:
\[Q_{\overline{\lambda}, m} = Q_{\lambda, m}\]

\begin{definition}
    Для дальнейшей работы нам понадобится ещё и такое пространство:
    \[Q_{\lambda, m}^k = \{x^k r(x) \; | \; r \in Q_{\lambda, m}\}, \; \; k \in \mathbb{Z}_{+}\]
\end{definition}

\begin{remark}
    Отметим следующие милые тривиальные свойства:
    \begin{enumerate}
        \item \[Q_{\lambda, m}^k \subset Q_{\lambda, m + k}\]
        \item \[Q_{\overline{\lambda}, m}^k = Q_{\lambda, m}^k\]
    \end{enumerate}
\end{remark}

\begin{remark}
    Если $\lambda = \alpha$, то есть $\beta = 0$, то
    \[Q_{\lambda, m} = \left\langle e^{\alpha x} x^t \; | \; t = 0, 1, \ldots, (m - 1)\right\rangle \]
    Если же $\lambda = \alpha + i \beta$, то есть $\beta \neq 0$, то 
    \[Q_{\lambda, m} = \left\langle e^{\alpha x} x^t \cos \beta x, e^{\alpha x} x^s \sin \beta x \; | \; t, s = 0, \ldots, (m - 1)\right\rangle \]
\end{remark}

\begin{exercise}
    Покажите, что все функции, входящие в набор
    \[\{e^{\alpha x} x^t \cos \beta x, e^{\alpha x} x^s \sin \beta x \; | \; t, s = 0, \ldots, (m - 1)\}\]
    являются линейно независимыми.
\end{exercise} \leavevmode
\\
Ещё для будущих выкладок нам пригодится такая форма записи введённых пространств:
\[\mu \in \mathbb{R} \implies Q_{\mu, m} = \left\langle \frac{e^{\mu x} x^{j - 1}}{(j - 1)!} \; | \; j = 1, \ldots, m\right\rangle \]

\[\mu = \alpha + i \beta, \; \beta \neq 0  \implies Q_{\mu, m} = \left\langle e^{\alpha x} \cdot \cos \beta x \cdot \frac{x^{j - 1}}{(j - 1)!}, \; e^{\alpha x} \cdot \sin \beta x \cdot \frac{x^{j - 1}}{(j - 1)!} \; | \; j = 1, \ldots, m\right\rangle \]

\begin{theorem}
    Пусть $F \in Q_{\mu, m}$, $\mu \in \mathbb{C}$, а $s$ --- кратность $\lambda = \mu$ как корня характеристического уравнения $\chi(\lambda) = 0$ ($s = 0$, если $\lambda = \mu$ не является корнем). Тогда у уравнения (\ref{eq:1}) есть частное решение, которое принадлежит $Q_{\mu, m}^s$.
\end{theorem}

\begin{proof}

    Проведем доказательство для $\mu \in \mathbb{R}$. Пусть $\mu$ --- корень $\chi(\lambda) = 0$ кратности $s \geq 0$. Тогда
    \[\chi(\lambda) = \zeta(\lambda)(\lambda - \mu)^s\]
    Здесь многочлен $\zeta(\mu) \neq 0$. Тогда уравнение (\ref{eq:1}) эквивалентно
    \begin{equation} \label{eq:3}
        \chi(\mathcal{D})(y(x)) = F(x)
    \end{equation}
    Но
    \[\chi(\mathcal{D}) = \zeta(\mathcal{D})(\mathcal{D} - \mu \mathcal{I})^s\]
    Покажем, что отображение $\chi(\mathcal{D}) \colon Q_{\mu, m}^s \to Q_{\mu, m}$ является биекцией. Тогда, обратив его, при помощи равенства (\ref{eq:3}) получим искомое частное решение. \\
    Достаточно показать, что каждое из отображений 
    \begin{equation*}
        (\mathcal{D} - \mu \mathcal{I})^s \colon Q_{\mu, m}^s \to Q_{\mu, m}
    \end{equation*}
    \begin{equation*}
        \zeta(\mathcal{D}) \colon Q_{\mu, m} \to Q_{\mu, m}
    \end{equation*}
    является биекцией в силу представления выше. \\
    Рассмотрим следующий базис в пространстве $Q_{\mu, m + s} \supset Q_{\mu, m}^s$:
    \[e_j(x) = \frac{e^{\mu x} x^{j - 1}}{(j - 1)!} \quad \forall j = 1, \ldots, (m + s)\]
    Понятно что оператор $\mathcal{D} \colon Q_{\mu, m + s} \to Q_{\mu, m + s}$. Кроме того, несложно видеть, что
    \[\mathcal{D} e_1 = \mu e_1 \quad \text{(} e_1 \text{ --- собственный)}\]
    \[\mathcal{D} e_j = \mu e_j + e_{j - 1}, \quad j = 2, \ldots, (m + s)\]
    Следовательно, матрица оператора $\mathcal{D}$ в этом базисе имеет вид
    \begin{equation*}
        J_{m + s}(\mu) = 
        \begin{pmatrix}
            \mu & 1 & 0 & 0 & 0 & \ldots & 0 \\
            0 & \mu & 1 & 0 & 0 & \ldots & 0 \\
            0 & 0 & \mu & 1 & 0 & \ldots & 0 \\
            \vdots & \vdots & \vdots & \vdots & \vdots & \vdots & \vdots \\
            0 & 0 & 0 & 0 & 0 & \ldots & 1 \\
            0 & 0 & 0 & 0 & 0 & \ldots & \mu
        \end{pmatrix}
        = \mu E_{m + s} + N_{m + s}, \text{ где}
    \end{equation*}
    \begin{equation*}
        E_{m + s} = 
        \begin{pmatrix}
            1 & 0 & 0 & 0 & 0 & \ldots & 0 \\
            0 & 1 & 0 & 0 & 0 & \ldots & 0 \\
            0 & 0 & 1 & 0 & 0 & \ldots & 0 \\
            \vdots & \vdots & \vdots & \vdots & \vdots & \vdots & \vdots \\
            0 & 0 & 0 & 0 & 0 & \ldots & 0 \\
            0 & 0 & 0 & 0 & 0 & \ldots & 1
        \end{pmatrix}
        , \quad N_{m + s} = 
        \begin{pmatrix}
            0 & 1 & 0 & 0 & 0 & \ldots & 0 \\
            0 & 0 & 1 & 0 & 0 & \ldots & 0 \\
            0 & 0 & 0 & 1 & 0 & \ldots & 0 \\
            \vdots & \vdots & \vdots & \vdots & \vdots & \vdots & \vdots \\
            0 & 0 & 0 & 0 & 0 & \ldots & 1 \\
            0 & 0 & 0 & 0 & 0 & \ldots & 0
        \end{pmatrix}
    \end{equation*}
    Матрица $N_{m + s}$ --- нильпотентная порядка $m + s$, то есть $(N_{m + s})^{m + s} = O$. \\
    Следовательно, у оператора $(\mathcal{D} - \mu \mathcal{I})^s$ матрица имеет вид $N_{m + s}^s$. Несложно понять, что
    \begin{equation*}
        N_{m + s}^s = 
        \left(
        \begin{array}{c|c}
        O_{m \times s} & E_m \\
        \hline
        O_{s \times s} & O_{s \times m}
        \end{array}
        \right)
    \end{equation*}
    Тогда матрица оператора $(\mathcal{D} - \mu \mathcal{I})^s$, ограниченная на подпространство $Q_{\mu, m}^s$, есть просто $E_m$. В самом деле, подпространство $Q_{\mu, m}^s$ получено из $Q_{\mu, m}$ домножением на моном $x^s$, а значит оно натянуто всё на те же векторы $e_j(x)$, только $j$ надо брать от $s$ до $(m + s)$. 
    Следовательно, $(\mathcal{D} - \mu \mathcal{I})^s$ --- биекция $Q_{\mu, m}^s$ и $Q_{\mu, m}$. \\
    Рассуждая для $\zeta(\mathcal{D})$ аналогично, просто рассмотрев базис в $Q_{\mu, m}$ вместо $Q_{\mu, m + s}$, получаем, что матрица
    оператора $\mathcal{D}$ в базисе $e_1, \ldots, e_m$ имеет вид
    \begin{equation*}
        J_m(\mu) = \mu E_m + N_m =
        \begin{pmatrix}
            \mu & 1 & 0 & 0 & 0 & \ldots & 0 \\
            0 & \mu & 1 & 0 & 0 & \ldots & 0 \\
            0 & 0 & \mu & 1 & 0 & \ldots & 0 \\
            \vdots & \vdots & \vdots & \vdots & \vdots & \vdots & \vdots \\
            0 & 0 & 0 & 0 & 0 & \ldots & 1 \\
            0 & 0 & 0 & 0 & 0 & \ldots & \mu
        \end{pmatrix}
    \end{equation*} 
    Несложно видеть, что
    \begin{equation*}
        J_m^l(\mu) = 
        \begin{pmatrix}
            \mu^l & \ast & \ast & \ast & \ast & \ldots & \ast \\
            0 & \mu^l & \ast & \ast & \ast & \ldots & \ast \\
            \vdots & \vdots & \vdots & \vdots & \vdots & \vdots & \vdots \\
            0 & 0 & 0 & 0 & 0 & \ldots & \ast \\
            0 & 0 & 0 & 0 & 0 & \ldots & \mu^l
        \end{pmatrix}
        , \text{ где } \ast \text{ --- некоторый элемент}
    \end{equation*}
    Как известно, матрица оператора $\zeta(\mathcal{D})$ получается из матрицы оператора $\mathcal{D}$ действием $\zeta$, то есть мы просто подставляем $J_m(\mu)$ в многочлен $\zeta$ и складываем соответствующие степени $J_m^l(\mu)$ с соответствующими коэффициентами. Несложно понять, в результате матрица $\zeta(\mathcal{D})$ будет иметь вид 
    \begin{equation*}
        \begin{pmatrix}
            \zeta(\mu) & \ast & \ast & \ast & \ast & \ldots & \ast \\
            0 & \zeta(\mu) & \ast & \ast & \ast & \ldots & \ast \\
            \vdots & \vdots & \vdots & \vdots & \vdots & \vdots & \vdots \\
            0 & 0 & 0 & 0 & 0 & \ldots & \ast \\
            0 & 0 & 0 & 0 & 0 & \ldots & \zeta(\mu)
        \end{pmatrix}
    \end{equation*}
    Но $\zeta(\mu) \neq 0$, а значит эта матрица невырожденна. Следовательно, $\zeta(\mathcal{D})$ --- биекция. Таким образом, требуемое доказано.
\end{proof}

\subsection{Метод вариации постоянных}

Пусть $y_1(x), \ldots, y_n(x)$ --- ФСР уравнения (\ref{eq:2}). Тогда по уже доказанному нами ранее, общее решение уравнения (\ref{eq:2}) имеет вид 
\[y(x) = \sum_{k = 1}^n c_k y_k(x), \text{ где } c_1, \ldots, c_n \in \mathbb{C}\]
Будем искать частное решение уравнения (\ref{eq:1}) в виде 
\begin{equation} \label{eq:4}
    y(x) = \sum_{k = 1}^n c_k(x) y_k(x)
\end{equation}

\begin{theorem}
    Пусть дифференцируемые функции $c_1, c_2, \ldots, c_n$ такие, что 
    \begin{equation*} \label{eq:ast} \tag{$\ast$}
        \begin{cases}
            \sum_{k = 1}^n c'_k(x) y_k(x) = 0 \\
            \sum_{k = 1}^n c'_k(x) y'_k(x) = 0 \\
            \ldots \\ 
            \sum_{k = 1}^n c'_k(x) y^{(n - 2)}_k(x) = 0 \\
            \sum_{k = 1}^n c'_k(x) y^{(n - 1)}_k(x) = \frac{F(x)}{a_0}
        \end{cases}
        \quad \forall x \in I \subset \mathbb{R}, \; \text{ где } I \text{ --- область определения } F
    \end{equation*}
    Тогда $y(x)$, заданная равенством (\ref{eq:4}), является решением уравнения (\ref{eq:1}).
\end{theorem}

\begin{proof}
    Для начала докажем лемму
    \begin{lemma}
        Пусть выполнена система $($\ref{eq:ast}$)$. Тогда
        \begin{equation} \label{eq:5}
            y^{(l)}(x) = \sum_{k = 1}^n c_k(x) y_k^{(l)}(x), \quad l = 0, 1, \ldots, (n - 1)
        \end{equation}
    \end{lemma}
    \begin{proof}
        Для $l = 0$ утверждение очевидно по определению (\ref{eq:4}) функции $y(x)$. При $l = 1$:
        \[y'(x) = \left(\sum_{k = 1}^n c_k(x) y_k(x)\right)' = \underbrace{\sum_{k = 1}^n c'_k(x) y_k(x)}_{0} + \sum_{k = 1}^n c_k(x) y'_k(x) = \sum_{k = 1}^n c_k(x) y'_k(x)\]
        Таким образом, получили базу индукции. \\
        Теперь проделаем шаг индукции. Пусть (\ref{eq:5}) выполнено для всех $l < n - 1$. Тогда
        \[y^{(l + 1)}(x) = \left(y^{(l)}(x)\right)' = \left(\sum_{k = 1}^n c_k(x) y_k^{(l)}(x)\right)' = \underbrace{\sum_{k = 1}^n c'_k(x) y_k^{(l)}(x)}_{0} + \sum_{k = 1}^n c_k(x) y_k^{(l + 1)}(x) = \sum_{k = 1}^n c_k(x) y_k^{(l + 1)}(x)\]
    \end{proof} \leavevmode
    \\
    Раз мы теперь знаем, как выглядят наши производные функции $y(x)$, просто возьмём и подставим их в уравнение (\ref{eq:1}):
    \[\sum_{j = 0}^n a_j y^{(n - j)}(x) = a_0 y^{(n)} (x) + \sum_{j = 1}^n a_j y^{(n - j)}(x) = a_0 \left( y^{(n - 1)} (x) \right)' + \sum_{j = 1}^n a_j y^{(n - j)}(x) = \]
    \[= a_0 \left(\sum_{k = 1}^n c_k(x) y_k^{(n - 1)}(x)\right)' + \sum_{j = 1}^n a_j y^{(n - j)}(x) = \]
    \[= a_0 \left(\sum_{k = 1}^n c'_k(x) y_k^{(n - 1)}(x) + \sum_{k = 1}^n c_k(x) y_k^{(n)}(x)\right) + \sum_{j = 1}^n a_j \sum_{k = 1}^n c_k(x) y_k^{(n - j)}(x) = \]
    \[= a_0 \cdot \frac{F(x)}{a_0} + \left(a_0 \sum_{k = 1}^n c_k(x) y_k^{(n)}(x) + \sum_{j = 1}^n a_j \sum_{k = 1}^n c_k(x) y_k^{(n - j)}(x) \right) = \]
    \[= F(x) + \sum_{j = 0}^n a_j \sum_{k = 1}^n c_k(x) y_k^{(n - j)}(x) = F(x) + \sum_{k = 1}^n c_k(x) \underbrace{\sum_{j = 0}^n a_j y_k^{(n - j)}(x)}_{0} = F(x)\]
    Следовательно, функция $y(x)$, заданная равенством (\ref{eq:4}), действительно суть решение (\ref{eq:1}), что и требовалось.
\end{proof}
\end{document}